\documentclass[11pt]{article}
\usepackage[utf8]{inputenc}
\usepackage{amsmath}
\usepackage{geometry}
\usepackage{listings}
\usepackage{xcolor}

\geometry{letterpaper, margin=1in}

\definecolor{codegreen}{rgb}{0,0.6,0}
\definecolor{backcolour}{rgb}{0.95,0.95,0.92}

\lstdefinestyle{mystyle}{
    backgroundcolor=\color{backcolour},   
    commentstyle=\color{codegreen},
    basicstyle=\ttfamily\footnotesize,
    breaklines=true,                 
    numbers=left,                    
    numbersep=5pt,                  
    tabsize=2
}

\lstset{style=mystyle}

\title{\textbf{Editorial: Problem C - Koyomi's Revisions}}
\author{Setter: Hudson}
\date{}

\begin{document}

\maketitle

\section*{Problem Overview}
Given a string $S$, answer $Q$ queries for the number of distinct substrings in $S[L \dots R]$.

\section*{Difficulty Analysis}
Rating: \textbf{2700}.
This is a well-known hard problem.
Naive solution (generating substrings and hashing) is $O(Q \cdot N^2)$, totally impossible.
Suffix Automaton construction is $O(N)$, but querying ranges is hard because the structure of the SAM corresponds to the \textit{whole} string.

\section*{Solution Approach: Offline Queries + SAM + Segment Tree}
We process queries **offline**, sorted by their right endpoint $R$.
As we iterate $R$ from $1$ to $N$, we want to maintain a data structure that allows us to sum up the contribution of valid substrings ending at or before $R$ that started at or after $L$.

A substring is defined by $(l, r)$. It contributes to the answer for query $[L, R]$ if $L \le l \le r \le R$.
To avoid overcounting duplicates (e.g., "a" appearing at index 1 and 3), we only count the **last occurrence** of each distinct substring.
If a distinct string $T$ appears multiple times ending at $r_1, r_2, \dots, r_k$, we only want to count the instance ending at $r_k$ (the rightmost one seen so far).

Algorithm:
1. Build a Suffix Automaton (SAM) for $S$.
2. We map each node $u$ in the SAM to the distinct substrings represented by that state. These substrings have lengths in range $(len(link(u)), len(u)]$.
3. We need to track the **last end position** of each SAM state. Let $pos(u)$ be the latest end-position of the occurrences of substrings in state $u$.
4. As we append character $S[R]$:
    - We traverse up the SAM link tree from the current state corresponding to the prefix $S[1 \dots R]$.
    - For each state $u$ on this path, the substrings represented by $u$ now have a new end position $R$.
    - Previously, they ended at $pos(u)$. We need to "remove" their contribution from $pos(u)$ and "add" it to $R$.
    - Specifically, state $u$ represents a range of lengths. The longest string in $u$ has length $len(u)$. The shortest has length $len(link(u)) + 1$.
    - This effectively updates a range of start positions.
5. This looks like "Color the path to root" in LCT.
    - We use a **Link-Cut Tree (LCT)** to maintain the SAM link tree.
    - The "Access" operation in LCT makes the path from the current node to the root a single "heavy" chain (or preferred path).
    - This corresponds to updating the `last\_pos` for all nodes on the path to $R$.
6. When the `last\_pos` of a chain changes from $old\_pos$ to $R$:
    - The substrings covered by this chain have lengths $[min\_len, max\_len]$.
    - Their starting positions were $[old\_pos - max\_len + 1, old\_pos - min\_len + 1]$.
    - We subtract 1 from these starting indices in a Segment Tree / Fenwick Tree.
    - We add 1 to the new starting indices $[R - max\_len + 1, R - min\_len + 1]$.
7. Answer for query $[L, R]$ is sum in Fenwick Tree range $[L, R]$.

Complexity: $O(N \log^2 N)$ using LCT + SegTree.

\end{document}
